\section{Introduction}
\label{sec:introduction}

Financial markets are often studied as complex systems in which heterogeneous agents, institutional constraints, liquidity cycles, information flows and feedback mechanisms interact across multiple time scales \cite{mantegna_stanley_1999_book,cont_2001_empirical}. Even when the analysis is restricted to daily asset returns, market variation is not fully described by the marginal behaviour of each series. It also has a relational component: how assets move together across the cross-section, summarized by their correlation matrix. This matrix reflects a mixture of collective market motion, sectoral organization, idiosyncratic shocks, crisis amplification and sampling noise. A central task is therefore to separate interpretable dependency structure from fluctuations that are compatible with random estimation effects \cite{raddant_dimatteo_2023}.

The distinction matters because financial risk is inherently multivariate. The volatility of an individual asset matters, but portfolio risk, contagion, diversification and systemic vulnerability depend on how assets co-move \cite{billio_2012_econometric}. During calm periods, correlations may appear moderate and dispersed. During stress periods, correlations often increase, reducing diversification benefits and lowering the effective dimensionality of the market \cite{junior_2012_correlation,zheng_2012_changes}.

Emerging equity markets are useful in this context because they combine global risk channels with local institutional, macroeconomic and liquidity conditions. Brazil is a suitable setting: B3 (Brasil, Bolsa, Balc\~ao), the country's stock exchange and financial-market infrastructure provider,\footnote{\url{https://www.b3.com.br/}} includes globally exposed commodity companies, large financial institutions, regulated utilities, consumer cyclicals, industrial firms, telecommunications firms and real-estate companies. These assets share exposure to domestic monetary policy, exchange-rate dynamics, country risk and global liquidity, but they also respond to distinct sectoral fundamentals. The dependency matrix therefore combines a dominant market component, group or sectoral components and a large residual component. Despite these features, the Brazilian equity market has received limited attention in the complex-systems literature. Previous studies have addressed Brazilian interbank networks \cite{tabak_2010_topology}, global financial system linkages \cite{silva_2016_structure}, agrarian cross-correlations \cite{stosic_2015_correlations}, multiscale networks \cite{pereira_2019_multiscale} and equity-market topology \cite{leonel_2022_brazilian}, but a comprehensive B3 equity study combining Random Matrix Theory filtering, planar financial networks and volatility forecasting remains insufficiently developed.

Correlation matrices are noisy, especially when the number of assets is not negligible relative to the number of observations. Random Matrix Theory (RMT) offers a benchmark for this problem by comparing the eigenspectrum with the spectrum expected from random correlation matrices \cite{marcenko_1967_distribution,laloux_1999_noise,plerou_2002_random}. Eigenvalues lying inside the Mar\v{c}enko--Pastur band are compatible with random sampling noise under the benchmark model, whereas eigenvalues outside the band indicate candidate collective modes that deserve further interpretation. Subsequent work has extended RMT-based methods for cleaning large covariance matrices \cite{bun_2017_cleaning}. We use RMT as an intermediate layer for reconstructing broad market, group-level and noise-compatible components before graph construction, not as a mechanical classifier of economic sectors.

Network filtering gives an additional representation of the same dependency system. By transforming correlations into distances, financial networks retain selected dependency links among assets and make it possible to inspect market backbones, local clusters and central nodes. The Minimum Spanning Tree (MST) emphasizes the sparsest connected structure \cite{mantegna_1999_hierarchical}, while the Planar Maximally Filtered Graph (PMFG) retains a denser planar representation with richer local connectivity \cite{tumminello_2005_tool,tumminello_2010_correlation}. We use these constructions to ask whether strong dependencies are organized by sectors, whether central nodes change after market-mode removal and whether the intermediate-scale organization of the market contains information beyond average correlation. Filtered networks are also relevant for risk analysis, where peripheral assets can offer better diversification properties than central hubs \cite{pozzi_2013_spread,aste_2010_correlation}.

We develop an integrated econophysics workflow for Brazilian equities traded on B3. The workflow has three stages. First, we document stylized facts, static correlations, sectoral correlation differences and rolling market synchronization. Second, we apply RMT to decompose the correlation matrix into market, group/sector and noise components, followed by hierarchical clustering and ordered heatmaps. Third, we convert original and filtered matrices into MST and PMFG networks, study hub reorganization and aggregate dependencies at the subsector level. A volatility forecasting experiment then tests whether the extracted features improve machine-learning and simple deep-learning forecasts under a strict temporal validation design.


The paper makes four contributions:
\begin{enumerate}
  \item We build a reproducible econophysics workflow for B3 equities using adjusted daily prices from 1998 to 2025, including return construction, stylized-fact checks, correlation analysis, sectoral comparison and rolling synchronization.
  \item We apply Random Matrix Theory to a 58-asset historical core universe and reconstruct market, group/sector and noise-compatible components of the empirical correlation matrix, yielding a spectral interpretation of Brazilian equity dependencies.
  \item We compare MST and PMFG networks built from original and RMT-filtered matrices, with emphasis on hub reorganization, clique structure, clustering, sectoral retention and subsector-level aggregation.
  \item We evaluate whether RMT- and network-derived features improve realized-volatility forecasting in machine-learning and simple deep-learning models, using a strict temporal split and nested feature-set design.
\end{enumerate}

The article is organized as follows. Section~\ref{sec:related_work} reviews financial dependencies, Random Matrix Theory, filtered financial networks and volatility forecasting. Section~\ref{sec:data} describes the data and asset universe. Section~\ref{sec:methodology} presents the methodology and the mathematical definitions used throughout the paper. Sections~\ref{sec:stylized_correlation}--\ref{sec:forecasting} report the results. Sections~\ref{sec:discussion} and~\ref{sec:limitations} discuss interpretation and limitations, and Section~\ref{sec:conclusion} presents the final considerations.
